% =================================================
% USC - EDIUS | HEADER COMPLETO PARA QUARTO
% Adaptado a las Recomendacións de maquetaxe da tese
% Universidad de Santiago de Compostela
% =================================================

% -------------------------------------------------
% FUENTE PRINCIPAL Y COMPATIBILIDAD XeLaTeX
% -------------------------------------------------
\usepackage{fontspec}
\setmainfont{Times New Roman}
\usepackage{microtype}
\usepackage[none]{hyphenat}

% -------------------------------------------------
% INTERLINEADO Y PÁRRAFOS
% -------------------------------------------------
\usepackage{setspace}
\singlespacing

\setlength{\parindent}{0.75cm}
\setlength{\parskip}{10pt}

% Evita estiramientos verticales artificiales.
\raggedbottom

% -------------------------------------------------
% ENCABEZADOS Y NUMERACIÓN
% -------------------------------------------------
\usepackage{fancyhdr}

\renewcommand{\chaptermark}[1]{\markboth{\thechapter.\ #1}{}}

\pagestyle{fancy}
\fancyhf{}

% Cambiar el nombre si se requiere otro autor/a.
\fancyhead[LE]{{\fontsize{10}{12}\selectfont\textsc{José Vicente Ordóñez Yaguache}}}
\fancyhead[RO]{{\fontsize{10}{12}\selectfont\nouppercase{\leftmark}}}

\fancyfoot[LE,RO]{\thepage}

\renewcommand{\headrulewidth}{0.4pt}
\renewcommand{\footrulewidth}{0pt}

% Inicio de capítulo: sin encabezado, con numeración visible en el pie exterior.
\fancypagestyle{plain}{%
  \fancyhf{}%
  \fancyfoot[LE,RO]{\thepage}%
  \renewcommand{\headrulewidth}{0pt}%
  \renewcommand{\footrulewidth}{0pt}%
}

% Páginas completamente en blanco sin encabezado ni numeración.
\usepackage{emptypage}

% -------------------------------------------------
% TÍTULOS Y EPÍGRAFES USC (Uso seguro con titlesec)
% -------------------------------------------------
\usepackage{titlesec}

% Capítulo
\titleformat{\chapter}[hang]
  {\normalfont\bfseries\fontsize{15}{18}\selectfont\raggedleft}
  {\thechapter.}
  {0.5em}
  {}

% Antes del título: 3 retornos de línea. Después: 1 retorno de línea.
\titlespacing*{\chapter}
  {0pt}
  {3\baselineskip}
  {1\baselineskip}

\titleformat{\part}[display]
  {\normalfont\bfseries\fontsize{15}{18}\selectfont\centering}
  {}
  {0pt}
  {\MakeUppercase{PARTE~\thepart.\space}}

% Nivel 1: 1.1 grosa, versaleta, sin tabular.
\titleformat{\section}
  {\normalfont\bfseries\scshape\fontsize{12}{14}\selectfont}
  {\thesection}
  {0.5em}
  {}

% Nivel 2: 1.1.1 grosa, tabulada a 0,75 cm.
\titleformat{\subsection}
  {\normalfont\bfseries\fontsize{12}{14}\selectfont}
  {\thesubsection}
  {0.75cm}
  {}

% Nivel 3: 1.1.1.1 normal, tabulada a 1,5 cm.
\titleformat{\subsubsection}
  {\normalfont\fontsize{12}{14}\selectfont}
  {\thesubsubsection}
  {1.5cm}
  {}

% Nivel 4: 1.1.1.1.1 cursiva, tabulada a 2,25 cm.
\titleformat{\paragraph}
  {\normalfont\itshape\fontsize{12}{14}\selectfont}
  {\theparagraph}
  {2.25cm}
  {}

% Espaciado entre epígrafes y texto.
\titlespacing*{\section}{0pt}{1\baselineskip}{1\baselineskip}
\titlespacing*{\subsection}{0pt}{1\baselineskip}{1\baselineskip}
\titlespacing*{\subsubsection}{0pt}{1\baselineskip}{1\baselineskip}
\titlespacing*{\paragraph}{0pt}{1\baselineskip}{0.5\baselineskip}

% Evitar títulos aislados al final de la página
\usepackage{needspace}

% -------------------------------------------------
% ÍNDICES
% -------------------------------------------------
\usepackage{tocloft}

\renewcommand{\contentsname}{ÍNDICE}
\renewcommand{\cfttoctitlefont}{\hfill\bfseries\fontsize{15}{18}\selectfont}
\renewcommand{\cftaftertoctitle}{\hfill}

\renewcommand{\listfigurename}{Índice de figuras}
\renewcommand{\listtablename}{Índice de tablas}

\renewcommand{\cftloftitlefont}{\hfill\bfseries\fontsize{15}{18}\selectfont}
\renewcommand{\cftafterloftitle}{\hfill}
\renewcommand{\cftlottitlefont}{\hfill\bfseries\fontsize{15}{18}\selectfont}
\renewcommand{\cftafterlottitle}{\hfill}

% Etiquetas de listas de figuras y tablas.
\renewcommand{\cftfigpresnum}{Figura~}
\renewcommand{\cftfigaftersnum}{: }
\setlength{\cftfignumwidth}{6em}

\renewcommand{\cfttabpresnum}{Tabla~}
\renewcommand{\cfttabaftersnum}{: }
\setlength{\cfttabnumwidth}{6em}

% -------------------------------------------------
% FIGURAS, LÁMINAS Y TABLAS
% -------------------------------------------------
\usepackage{caption}

\newfontfamily\trebuchet{Trebuchet MS}
\DeclareCaptionFont{uscCaptionFont}{\trebuchet\bfseries\fontsize{9}{11}\selectfont}

\captionsetup{
  font=uscCaptionFont,
  labelfont=bf,
  textfont=bf,
  justification=centering,
  singlelinecheck=false,
  skip=6pt
}

% Texto de tablas en Trebuchet MS 9 pt.
\usepackage{etoolbox}

\def\uscNoShortCaption{\uscNoShortCaption}
\long\def\uscEmitFigureCaption#1#2{%
  \def\uscShortCaptionValue{#1}%
  \ifx\uscShortCaptionValue\uscNoShortCaption
    \uscOriginalCaption{#2}%
  \else
    \uscOriginalCaption[#1]{#2}%
  \fi
}
\long\def\uscInsertFigureSourceBreak#1#2 Fuente:#3 Fuente:#4\uscSourceEnd{%
  \ifblank{#4}{%
    \uscEmitFigureCaption{#1}{#2}%
  }{%
    \uscEmitFigureCaption{#1}{#2\newline Fuente:#3}%
  }%
}
\long\def\uscPreserveFigureSourceBreak#1#2\newline Fuente:#3\newline Fuente:#4\uscSourceEnd{%
  \ifblank{#4}{%
    \uscInsertFigureSourceBreak{#1}#2 Fuente:\uscSourceMarker{} Fuente:\uscSourceEnd
  }{%
    \uscEmitFigureCaption{#1}{#2\newline Fuente:#3}%
  }%
}
\long\def\uscProcessFigureCaption#1#2{%
  \uscPreserveFigureSourceBreak{#1}#2%
    \newline Fuente:\uscSourceMarker\newline Fuente:\uscSourceEnd
}
\long\def\uscFigureCaptionWithoutShort#1{%
  \uscProcessFigureCaption{\uscNoShortCaption}{#1}%
}
\long\def\uscFigureCaptionWithShort[#1]#2{%
  \uscProcessFigureCaption{#1}{#2}%
}
\makeatletter
\AtBeginEnvironment{figure}{%
  \let\uscOriginalCaption\caption
  \renewcommand{\caption}{%
    \@ifnextchar[{\uscFigureCaptionWithShort}{\uscFigureCaptionWithoutShort}%
  }%
}
\makeatother

\AtBeginEnvironment{tabular}{\trebuchet\fontsize{9}{11}\selectfont}
\AtBeginEnvironment{longtable}{\trebuchet\fontsize{9}{11}\selectfont}

\usepackage{threeparttable}
\usepackage{booktabs}
\usepackage{longtable}
\usepackage{array}

% -------------------------------------------------
% NOTAS AL PIE
% -------------------------------------------------
\usepackage[hang,flushmargin]{footmisc}
\renewcommand{\footnotesize}{\fontsize{9}{11}\selectfont}

% -------------------------------------------------
% CITAS TEXTUALES LARGAS
% -------------------------------------------------
\usepackage{changepage}

\newenvironment{longquote}
{%
  \par\vspace{\baselineskip}%
  \begin{adjustwidth}{2.25cm}{1.5cm}%
  \fontsize{11}{13}\selectfont%
  \setlength{\parindent}{0pt}%
  \justifying%
}
{%
  \end{adjustwidth}%
  \vspace{\baselineskip}%
}

% -------------------------------------------------
% JUSTIFICACIÓN Y CONTROL DE LÍNEAS
% -------------------------------------------------
\usepackage{ragged2e}
\justifying

\clubpenalty=10000
\widowpenalty=10000
\displaywidowpenalty=10000
\brokenpenalty=10000

% Desactiva globalmente la separación silábica automática.
\hyphenpenalty=10000
\exhyphenpenalty=10000
\tolerance=2500
\emergencystretch=3em

% -------------------------------------------------
% LISTAS
% -------------------------------------------------
\usepackage{enumitem}

\setlist[itemize]{
  topsep=3pt,
  partopsep=0pt,
  parsep=0pt,
  itemsep=3pt,
  leftmargin=1.25cm
}

\setlist[enumerate]{
  topsep=3pt,
  partopsep=0pt,
  parsep=0pt,
  itemsep=3pt,
  leftmargin=1.25cm
}

% -------------------------------------------------
% FLOTANTES
% -------------------------------------------------
\usepackage{graphicx}
\usepackage{float}
\usepackage{placeins}
\usepackage{flafter}

\setcounter{topnumber}{3}
\setcounter{bottomnumber}{2}
\setcounter{totalnumber}{5}
\renewcommand{\topfraction}{0.9}
\renewcommand{\bottomfraction}{0.8}
\renewcommand{\textfraction}{0.07}
\renewcommand{\floatpagefraction}{0.75}

% Impide que un flotante pendiente invada la sección siguiente.
% Se conserva la colocación natural de LaTeX dentro de cada apartado.
\pretocmd{\section}{\FloatBarrier}{}{}

% -------------------------------------------------
% APÉNDICES
% -------------------------------------------------
\usepackage[toc,page]{appendix}
\renewcommand{\appendixname}{Anexo}
\renewcommand{\appendixtocname}{ANEXOS}
\renewcommand{\appendixpagename}{ANEXOS}

% -------------------------------------------------
% HIPERVÍNCULOS
% -------------------------------------------------
\PassOptionsToPackage{hidelinks}{hyperref}

% -------------------------------------------------
% COMANDOS ÚTILES PARA PÁGINAS PRELIMINARES
% -------------------------------------------------
\newcommand{\USCPrelimTitle}[1]{%
  \begin{center}
    {\bfseries\fontsize{15}{18}\selectfont #1}
  \end{center}
  \vspace{1\baselineskip}
}

\newcommand{\USCAbbreviationsTitle}{\USCPrelimTitle{LISTA DE ABREVIATURAS}}
\newcommand{\USCPublicationsTitle}{\USCPrelimTitle{LISTA DE PUBLICACIONES DERIVADAS DE LA TESIS}}

% =================================================
